\documentclass{article}

% NeurIPS 2025 style (vendored in neurips/neurips_2025.sty), preprint mode:
% this is a self-study report, not a conference submission.
\PassOptionsToPackage{numbers,sort&compress}{natbib}
\usepackage[preprint]{neurips/neurips_2025}

\usepackage[T1]{fontenc}
\usepackage{amsmath,amssymb}
\usepackage{booktabs}
\usepackage{graphicx}
\usepackage[colorlinks=true,linkcolor=blue,urlcolor=blue,citecolor=blue]{hyperref}

\title{Why Scaled Dot-Product Attention Divides by $\sqrt{d_k}$}
\author{self-study-with-ai}
\date{\today}

\begin{document}
\maketitle

\begin{abstract}
The Transformer computes attention weights with softmax over query-key dot
products divided by $\sqrt{d_k}$. Vaswani et~al.~\citep{vaswani2017attention}
state their suspicion plainly: at large key dimensionality, unscaled dot
products grow large in magnitude and push softmax into regions of extremely
small gradients. This report reconstructs that argument from the primary
source and marks where the source asserts rather than demonstrates.
\end{abstract}

\section{Introduction}

The brief asks why the Transformer scales attention logits by
$1/\sqrt{d_k}$. The answer, per the primary source, is a gradient-magnitude
argument about the softmax function, offered as the mechanism that makes the
computationally cheaper dot-product attention usable at large $d_k$.

\section{Background}

Attention maps a query and key--value pairs to an output as a weighted sum of
the values; the weight on each value depends on a compatibility function of
the query with the corresponding key~\citep{vaswani2017attention}. Two
compatibility functions are common: additive attention, which computes the
score with a single-hidden-layer feedforward network, and dot-product
attention, which uses the raw inner product~\citep{vaswani2017attention}.

\section{Methods and Sources}

Scope: a single-source briefing on Vaswani et~al.~\citep{vaswani2017attention}
(NeurIPS 2017), the paper that introduced the mechanism. Coverage is limited to what
that paper states about the scaling factor; attention variants and later
empirical work are out of scope.

\section{Findings}

The Transformer uses scaled dot-product attention~\citep{vaswani2017attention}:
\begin{equation}
\mathrm{Attention}(Q, K, V) = \mathrm{softmax}\!\left(\frac{QK^\top}{\sqrt{d_k}}\right) V ,
\end{equation}
where $Q$, $K$, $V$ are matrices of queries, keys, and values and $d_k$ is the
key dimensionality. The authors' stated motivation~\citep{vaswani2017attention}:
for large $d_k$, the dot products grow large in magnitude, pushing the
softmax function into regions where it has extremely small gradients;
dividing by $\sqrt{d_k}$ counteracts this effect.

The engineering context matters. The authors note that dot-product attention
is faster and more space-efficient than additive attention because it reduces
to highly optimized matrix multiplication, and that the two forms are
theoretically similar in complexity; the scaling factor is what keeps
dot-product attention competitive with additive attention at large
$d_k$~\citep{vaswani2017attention}.

What the argument claims: if each logit is a sum of $d_k$ products of roughly
unit-scale quantities, its magnitude grows with $d_k$; dividing by
$\sqrt{d_k}$ corresponds to normalizing by the standard deviation of the sum
under a unit-variance assumption, keeping logits in the region where softmax
gradients remain usable. This middle step is a standard reading of the
paper's stated motivation, not a sentence from the paper itself.

In multi-head attention, the mechanism operates with $h = 8$ heads and
$d_k = d_v = d_{\mathrm{model}} / h = 64$ at $d_{\mathrm{model}} =
512$~\citep{vaswani2017attention}.

\section{Related Work}

Out of scope for this briefing: additive attention~\citep{vaswani2017attention}
is the comparison point named by the source; later refinements and
alternatives to the softmax are not covered.

\section{Limitations}

The scaling motivation is asserted, not measured: Vaswani
et~al.~\citep{vaswani2017attention} provide neither an ablation of scaled
versus unscaled dot-product attention nor numeric support for the claim that
unscaled dot-product attention underperforms additive attention at large
$d_k$. The standard-deviation reading in the Findings is an interpretation,
labeled as such. Coverage is one primary source; later empirical work on the
claim is outside the search scope.

\section{Conclusion}

The division by $\sqrt{d_k}$ exists to keep attention logits in the range
where softmax gradients stay usable as key dimensionality grows, which is
what makes the cheaper dot-product form viable in the Transformer. The
justification in the primary source is analytical and stated with hedging;
confirming it empirically requires additional sources outside this briefing.

\bibliographystyle{plainnat}
\bibliography{refs}

\end{document}
